\section{实验任务与场地}
\subsection{任务要求与验收口径}
依据课程《第五--六章：基于视觉与 AI 定位的多机协同控制实验》\cite{course}，系统同时包含人形机器人搬运和多车跟随两项任务。TonyPi 从指定位置抱起目标物体，在障碍通道中至少完成两个转角（左右转各至少一次），到达终点后放下物体；TurtleBot2 保持跟随并在终点区域停靠。课程总体评价由现场竞赛与方案、源码报告构成，权重分别为80\%和20\%。手册中的任务平台数量写为3--4台 TurtleBot2，详细评分表按四台、四种方式设置；本报告为四种方式保留独立章节，实际参赛车辆与方式认定须由全队和教师确认。

\begin{table}[htbp]\centering\small
\caption{任务要求与对应记录}
\begin{tabularx}{\linewidth}{p{30mm}YY}\toprule
要求 & 系统对应内容 & 应提交的证据\\\midrule
抱持、转弯与放置 & 抓取动作、上肢锁定、A左转/B右转/C停止 & 连续视频、动作文件版本、决策日志\\
不接触障碍物 & 人形路线约束；TurtleBot代价地图与门控 & 接触次数、障碍位置、最小间距实测\\
保持跟随目标 & 标记身份、目标时效检查、失效处置 & 丢失次数、持续时间、重捕获记录\\
多种独立方式 & 四组分别说明观测来源和控制链 & 各组实现、接口及现场运行证据\\
协同到达终点 & 任务状态、各车停靠条件 & 全队同场运行记录与到达时间\\\bottomrule
\end{tabularx}\end{table}

\subsection{场地与感知约束}
现场使用灰色地板及相对连续的灰色挡板构成通道，路线包含直行、左转、直行、右转和终点区域。地板接缝、胶带及其他地面标识均视为无关背景，不能作为路线识别标准。目标物块由早期蓝色调整为红色；抱持后物块可能遮挡画面下部，因此抓取确认与前方导航需要分别配置观察区域或头部姿态。

黄色雪糕筒可作为早期道路居中方案的终点候选；最新标记路线以固定标记 C 的距离触发停止。两种终点条件属于不同程序配置，不应在同一次实验记录中混写。图\ref{fig:route}仅表达路径拓扑，方向和长度以现场测量为准。
\begin{figure}[htbp]\centering
\begin{tikzpicture}[x=1cm,y=1cm]
\draw[line width=10mm,gray!15] (0,0)--(4,0)--(4,2)--(8,2);
\draw[flow,blue!70!black] (0,0)--(4,0)--(4,2)--(8,2);
\foreach \p/\t in {(0,0)/起点,(4,0)/A：左转,(4,2)/B：右转,(8,2)/C：停止放置}
 {\fill \p circle(2pt);\node[above=5mm,font=\small] at \p {\t};}
\node[below,font=\small] at (2,0) {第一段直行};
\node[right,font=\small] at (4,1) {第二段};
\node[below,font=\small] at (6,2) {第三段直行};
\end{tikzpicture}

\caption{路线拓扑示意；不按比例绘制，地面标识不参与定位}\label{fig:route}
\end{figure}
\reportphoto{figures/site/corridor.jpg}{灰色挡板通道现场}{选用通道全景；注明拍摄位置、日期和行进方向。}
